% Quiz: Prompt Engineering (Weeks 6 & 8) - Multiple Choice (12 questions)
\documentclass[11pt,a4paper]{article}
\usepackage[utf8]{inputenc}
\usepackage{geometry}
\usepackage{enumitem}
\geometry{margin=1in}
\begin{document}
\begin{center}
  {\LARGE Prompt Engineering \textemdash{} Quiz 4}\\[6pt]
  {\large Name: \underline{\hspace{6cm}} \quad Date: \underline{\hspace{3cm}}}
\end{center}
\vspace{0.5cm}

\textbf{Instructions:} This quiz contains 12 multiple-choice questions based on the content from Iterative Refinement and Evaluation. Answer all questions.
\vspace{0.5cm}
\begin{enumerate}[leftmargin=*, label=\textbf{\#\arabic*}]

  \item Which error category describes a situation where the model ignores a negative constraint (e.g., "Do not use lists")?\\
    \begin{enumerate}[label=\Alph*)]
      \item Reasoning Failure
      \item Constraint Violation
      \item Tone Drift
      \item Phrasing Sensitivity
    \end{enumerate}

  \item In the "Critique Loop" pattern, what is the primary purpose of the critique step?\\
    \begin{enumerate}[label=\Alph*)]
      \item To generate the final answer immediately.
      \item To ask the model to evaluate its own output against a rubric before revising.
      \item To ask a human to rewrite the prompt.
      \item To increase the temperature of the model.
    \end{enumerate}

  \item What is the main risk associated with "Vibe Coding" (iterating based on "look and feel" rather than code inspection)?\\
    \begin{enumerate}[label=\Alph*)]
      \item It is too slow for prototyping.
      \item It requires deep knowledge of Python.
      \item Fragility and inability to fix edge cases because the user doesn't understand the code.
      \item It produces code that is too verbose.
    \end{enumerate}

  \item What is "LLM-as-a-Judge"?\\
    \begin{enumerate}[label=\Alph*)]
      \item Using a lawyer to review AI outputs.
      \item Using a strong model (e.g., GPT-4) to grade the outputs of other models.
      \item A technique to fine-tune models on legal data.
      \item A method to generate court transcripts.
    \end{enumerate}

  \item Which phase of the Prompt Engineering Lifecycle involves checking outputs against requirements like accuracy and tone?\\
    \begin{enumerate}[label=\Alph*)]
      \item Draft
      \item Evaluate
      \item Refine
      \item Deploy
    \end{enumerate}

  \item The "Evaluation Gap" refers to the paradox that:\\
    \begin{enumerate}[label=\Alph*)]
      \item It is harder to generate text than to verify it.
      \item It is significantly easier to generate text than to verify it.
      \item Models are too slow to evaluate.
      \item Humans are faster at generating text than AI.
    \end{enumerate}

  \item Which of the following is NOT one of the "3 H's" of evaluation?\\
    \begin{enumerate}[label=\Alph*)]
      \item Helpful
      \item Honest
      \item Harmless
      \item Human-like
    \end{enumerate}

  \item What is a "Golden Dataset"?\\
    \begin{enumerate}[label=\Alph*)]
      \item A large dataset used for pre-training the model.
      \item A curated set of inputs and "perfect" outputs used as a benchmark for testing.
      \item The most expensive dataset available for purchase.
      \item A set of rejected prompts.
    \end{enumerate}

  \item Which of the following is considered a "Hard Metric"?\\
    \begin{enumerate}[label=\Alph*)]
      \item Creativity
      \item Empathy
      \item Syntax or JSON format compliance
      \item Flow
    \end{enumerate}

  \item "Verbosity Bias" refers to the tendency of models (and humans) to:\\
    \begin{enumerate}[label=\Alph*)]
      \item Prefer shorter, more concise answers.
      \item Rate longer answers as better, even if they contain fluff.
      \item Prefer answers that appear first in a list.
      \item Prefer answers generated by themselves.
    \end{enumerate}

  \item Which sampling strategy involves testing high-stakes domains or new templates more frequently?\\
    \begin{enumerate}[label=\Alph*)]
      \item Random sampling
      \item Risk-based sampling
      \item Stratified sampling
      \item Event-driven sampling
    \end{enumerate}

  \item In a "Pairwise Comparison" evaluation, the judge is asked to:\\
    \begin{enumerate}[label=\Alph*)]
      \item Rate a single response from 1 to 5.
      \item Compare the response to a Golden Answer.
      \item Decide which of two responses (A or B) is better.
      \item Check for forbidden words.
    \end{enumerate}

\end{enumerate}
\end{document}
